\documentclass[10pt,a4paper]{amsart}
\usepackage[margin=19mm]{geometry}
\usepackage{amsmath,amssymb,mathtools}
\usepackage[scr=rsfs,cal=euler]{mathalfa}
\usepackage{xfrac}
\usepackage[unicode,hidelinks,bookmarks=false]{hyperref}
\newcommand{\ie}{i.e.\ }
\newcommand{\cf}{cf.\ }
\newcommand{\resp}{respectively\ }
\newcommand{\Subsection}{\subsection}
\providecommand{\nobreakdash}{\nobreak\hbox{-}}
\setlength{\emergencystretch}{2em}
\multlinegap0pt
\usepackage{bm}
\usepackage{bbm}
\usepackage{dsfont}
\usepackage{xspace}
\usepackage{cancel}
\usepackage{mathtools}
\usepackage{amsthm}
\newtheorem{theorem}{Theorem}
\title[Blind error estimation for CUR approximation]{Blind error estimation for CUR approximation (Theorem 3)}
\author{Lorenzo Lazzarino, Katherine J. Pearce,, Nathaniel Pritchard}
\date{Source: arXiv:2609.23810v1 (CC BY 4.0)}
\begin{document}
\maketitle

\noindent\emph{Source and licence.} Blind error estimation for CUR approximation. Version arXiv:2609.23810v1 (CC BY 4.0), distributed under the Creative Commons Attribution 4.0 International licence, as recorded in the arXiv licence metadata for this exact version and shown on the abstract page https://arxiv.org/abs/2609.23810v1. Full rights notice: licenses/arxiv-2609.23810-CC-BY-4.0.txt.\par\smallskip

\noindent\textit{Editorial scope.} Counted unit: the Theorem printed as Theorem 3 of the article (source0.tex lines 943--1010, source label \texttt{None}), delivered together with the article's own complete original proof of it (source0.tex lines 1012--1155, one \texttt{proof} environment of 144 lines). No mathematics is written, repaired or re-derived: statement and proof are the article's own text line for line. Required context retained uncounted: none. Every definition, display and label of those lines is retained unchanged. Omitted from this package and not redistributed: 1--942, 1011--1011, 1156--1933. Nothing inside a retained excerpt is omitted or altered: every retained body is delivered exactly as the article's lines are written, so no omission receipt of any kind accompanies this package. Citations: the retained excerpts carry no citation command at all, so no bibliography entry of the article is transcribed and no reference is redistributed. Reference adaptation: none was needed and none was made; every reference inside the delivered unit resolves inside it, so no unresolved reference or citation is produced. Printed slips kept literal: no printed word, symbol, hypothesis or number of the counted statement or of its proof was altered. Layout and class: the article's own class is \texttt{article} with packages including geometry, natbib, graphicx, epstopdf, pgfplots, pgfplotstable, amsmath, amssymb, amsfonts, amscd, amsthm, xy, enumerate, enumitem; the wrapper is the lane's pinned \texttt{amsart} editorial wrapper at 10pt on A4, which restates the theorem-like environments the retained text needs and adds the article's own macro definitions that the retained text needs (none), with extra wrapper packages (bm, bbm, dsfont, xspace, cancel, mathtools). The delivered numbering is the unit's own. Assets: no retained span contains a figure environment, a graphics-inclusion command, a PDF-page inclusion, a table or a TikZ picture; no image file is copied into this package, the package declares no asset dependency and no image file is redistributed with it. Licence: version arXiv:2609.23810v1 of this article is distributed under the Creative Commons Attribution 4.0 International licence, as recorded in the arXiv licence metadata for this exact version and shown on the abstract page https://arxiv.org/abs/2609.23810v1. A full rights notice is delivered at licenses/arxiv-2609.23810-CC-BY-4.0.txt. Characters: every retained excerpt is pure ASCII, so the delivered engine, XeLaTeX with its default Latin Modern text font, reports no missing character.
\begin{theorem}
    Consider the family of matrices
    \begin{equation}
        \mathbf{A}^{(0)}
        =
        \begin{pmatrix}
            \mathbf{L} & \mathbf{0} \\
            0          & \mathbf{0}
        \end{pmatrix}
        \in\mathbb{R}^{m\times n},
        \qquad
        \mathbf{A}^{(j)}
        =
        \begin{pmatrix}
            \mathbf{L} & \mathbf{0}         \\
            0          & h\mathbf{e}_j^\top
        \end{pmatrix}
        \in\mathbb{R}^{m\times n},
        \qquad j=1,\ldots,d,
    \end{equation}
    where $\mathbf{L}\in\mathbb{R}^{\bar m\times\bar n}$ has rank one and
    $\|\mathbf{L}\|_F=1$, $m=\bar m+1$, $n=\bar n+d$, $h>0$, and
    $\mathbf{e}_j\in\mathbb{R}^d$ denotes the $j$-th canonical basis vector.
    Assume a CUR approximation is given, and that, for each $\mathbf{A}^{(j)}$,
    the column corresponding to $h$ is not captured by the approximation, i.e.,
    the CUR residuals are
    \begin{equation}
        \mathbf{A}^{(0)}-\mathbf{A}^{(0)}_{\mathrm{CUR}}
        =
        \begin{pmatrix}
            \mathbf{F}_L & \mathbf{0} \\
            0            & \mathbf{0}
        \end{pmatrix},
        \qquad
        \mathbf{A}^{(j)}-\mathbf{A}^{(j)}_{\mathrm{CUR}}
        =
        \begin{pmatrix}
            \mathbf{F}_L & \mathbf{0}         \\
            0            & h\mathbf{e}_j^\top
        \end{pmatrix},
        \qquad j=1,\ldots,d,
    \end{equation}
    where $\|\mathbf{F}_L\|_F=\varepsilon_L$. Define
    \begin{equation}
        \mathcal{T}
        :=
        \left\{
        \bar n+j:\ j\in\{1,\ldots,d\}
        \right\},
    \end{equation}
    and fix $\delta>1$. Assume that
    $h^2>(\delta^2-1)\varepsilon_L^2$. Finally, assume that all information available to the estimator prior to its additional column queries is identical across the matrices $\mathbf{A}^{(0)}, \dots,\mathbf{A}^{(d)}$

    Then, for every target success probability $\pi\in(1/2,1)$, there does not
    exist a randomized, possibly adaptive, algorithm that inspects at most
    \begin{equation}
        q < (2\pi-1)|\mathcal{T}|
    \end{equation}
    columns and outputs a $\delta$-approximation of
    $\|\mathbf{A}-\mathbf{A}_{\mathrm{CUR}}\|_F^2$ with probability at least
    $\pi$ for every input in
    $\{\mathbf{A}^{(0)},\mathbf{A}^{(1)},\ldots,\mathbf{A}^{(d)}\}$.
    In particular, any such algorithm requires
    $$
        q=\Omega(|\mathcal{T}|)=\Omega(d)
    $$
    column queries in the worst case.
\end{theorem}

\begin{proof}
    Consider the distribution $\mathcal{D}$ over the input matrices defined as
    follows. With probability $1/2$, draw $\mathbf{A}^{(0)}$, while with
    probability $1/2$, draw $\mathbf{A}^{(J)}$, where
    \begin{equation}
        J\sim\operatorname{Unif}\{1,\ldots,d\}.
    \end{equation}
    Equivalently,
    \begin{equation}
        \mathbb{P}_{\mathbf{A}\sim\mathcal{D}}
        \left(\mathbf{A}=\mathbf{A}^{(0)}\right)
        =\frac{1}{2},
        \qquad
        \mathbb{P}_{\mathbf{A}\sim\mathcal{D}}
        \left(\mathbf{A}=\mathbf{A}^{(j)}\right)
        =\frac{1}{2d},
        \qquad j=1,\ldots,d.
    \end{equation}

    Fix a deterministic algorithm that adaptively inspects at most $q$ columns.
    Run the algorithm on the input $\mathbf{A}^{(0)}$, and collect the indices
    of the inspected columns in $\mathcal{J}_q$. Since the algorithm is
    deterministic, $\mathcal{J}_q$ is fixed by the observations made on
    $\mathbf{A}^{(0)}$. Let
    \begin{equation}
        \mathcal{S}
        :=
        \mathcal{J}_q\cap\mathcal{T}
    \end{equation}
    denote the set of inspected columns among the $d$ possible locations of the
    additional component. Clearly,
    \begin{equation}
        |\mathcal{S}|\leq q.
    \end{equation}

    Now consider an input $\mathbf{A}^{(j)}$ with
    $\bar n+j\notin\mathcal{S}$. The only difference between
    $\mathbf{A}^{(0)}-\mathbf{A}^{(0)}_{\mathrm{CUR}}$ and
    $\mathbf{A}^{(j)}-\mathbf{A}^{(j)}_{\mathrm{CUR}}$ lies in column
    $\bar n+j$. Since this column is not inspected along the execution of the
    algorithm on $\mathbf{A}^{(0)}$, every column observed by the algorithm is
    identical under $\mathbf{A}^{(0)}$ and $\mathbf{A}^{(j)}$. Therefore, by
    induction over the adaptive queries, the deterministic algorithm makes the
    same sequence of queries and observes the same values on both inputs.
    Consequently, it must output the same value for $\mathbf{A}^{(0)}$ and
    $\mathbf{A}^{(j)}$.

    The target quantities are
    \begin{equation}
        \left\|
        \mathbf{A}^{(0)}-\mathbf{A}^{(0)}_{\mathrm{CUR}}
        \right\|_F^2
        =
        \varepsilon_L^2,
        \qquad
        \left\|
        \mathbf{A}^{(j)}-\mathbf{A}^{(j)}_{\mathrm{CUR}}
        \right\|_F^2
        =
        \varepsilon_L^2+h^2.
    \end{equation}
    Since
    $h^2>(\delta^2-1)\varepsilon_L^2$,
    we have
    \begin{equation}
        \delta\varepsilon_L^2
        <
        \frac{\varepsilon_L^2+h^2}{\delta}.
    \end{equation}
    Hence the two $\delta$-approximation intervals
    \begin{equation}
        \mathcal{I}_0
        =
        \left[
            \frac{\varepsilon_L^2}{\delta},
            \delta\varepsilon_L^2
            \right],
        \qquad
        \mathcal{I}_j
        =
        \left[
            \frac{\varepsilon_L^2+h^2}{\delta},
            \delta(\varepsilon_L^2+h^2)
            \right]
    \end{equation}
    are disjoint. Therefore, whenever $\bar n+j\notin\mathcal{S}$, the common
    output of the deterministic algorithm on $\mathbf{A}^{(0)}$ and
    $\mathbf{A}^{(j)}$ can be a $\delta$-approximation for at most one of these
    two inputs.

    We now bound the average success probability of the deterministic algorithm
    under $\mathcal{D}$. Let $a_0\in\{0,1\}$ indicate whether the algorithm is
    successful on $\mathbf{A}^{(0)}$. For each $j$ such that
    $\bar n+j\notin\mathcal{S}$, the disjointness of the two approximation
    intervals implies that, if $a_0=1$, the algorithm must fail on
    $\mathbf{A}^{(j)}$. For indices satisfying $\bar n+j\in\mathcal{S}$, we
    upper bound the success probability by one. Therefore, if the algorithm is
    successful on $\mathbf{A}^{(0)}$, its success probability under
    $\mathcal{D}$ is at most
    \begin{equation}
        \frac{1}{2}
        +
        \frac{1}{2}\frac{|\mathcal{S}|}{d}
        \leq
        \frac{1}{2}
        +
        \frac{q}{2d}.
    \end{equation}
    If the algorithm is unsuccessful on $\mathbf{A}^{(0)}$, then even assuming
    that it succeeds on every $\mathbf{A}^{(j)}$, its success probability under
    $\mathcal{D}$ is at most $1/2$. Thus, in either case,
    \begin{equation}
        \mathbb{P}_{\mathbf{A}\sim\mathcal{D}}
        \left(
        \operatorname{alg}\text{ successfully produces a }
        \delta\text{-approximation}
        \right)
        \leq
        \frac{1}{2}
        +
        \frac{q}{2d}.
    \end{equation}

    It follows that if
    \begin{equation}
        q < (2\pi-1)d
        =
        (2\pi-1)|\mathcal{T}|,
    \end{equation}
    then every deterministic algorithm using at most $q$ column queries has
    average success probability strictly smaller than $\pi$ under
    $\mathcal{D}$. By Yao's minimax principle, the same lower bound applies to
    randomized algorithms: there does not exist a randomized, possibly adaptive,
    algorithm using at most $q$ column queries that achieves success probability
    at least $\pi$ on every input in the support of $\mathcal{D}$. Since
    $$
        \operatorname{supp}(\mathcal{D})
        \subseteq
        \{
        \mathbf{A}^{(0)},\mathbf{A}^{(1)},\ldots,\mathbf{A}^{(d)}
        \},
    $$
    this proves the claim.
\end{proof}

\end{document}
